\begin{proof}

\begin{enumerate}[label={$\langle1\rangle$\arabic*.},leftmargin=*]
\item The integral Leray sequence gives an injection
\[
H^0(\PP^1,R^1f_*\Z)\longrightarrow H^2(\PP^1,\Z)=\Z
\]
whose cokernel injects into the finite group $H^2(Y,\Z)$. Thus its source is infinite cyclic. The degree-one specialization argument of Part D identifies its stalk at every fibre with a nonzero integral class: a class in the preimage of a small disc that vanishes on a nearby fibre comes from the meridian of the punctured disc; a normal disc to any irreducible component makes a nonzero multiple of that meridian nullhomotopic, forcing the class to vanish. The real-to-holomorphic cohomology map on a compact connected complex space is injective on real classes, as proved there by the pluriharmonic maximum principle. Coherent Leray therefore supplies $\beta\in H^1(Y,\OO_Y)$ whose restriction to every full fibre is the image of $\gamma$.

\item Fix a real irrational number $r$ and put $M_r=\exp(r\beta)$, with the exponential normalized by $2\pi i$. On each fibre this line is represented by the unitary character $\exp(2\pi ir\gamma)$. It is nontrivial, and on a singular fibre its normalization pullback is nontrivial by the hypothesis on $\gamma$. To see this directly, triviality of an $n$th tensor power would place $nr\gamma$ in the integral lattice after applying $H^1(\R)\to H^1(\OO)$; real injectivity makes this an equality in $H^1(\R)$, which is impossible for an irrational multiple of a nonzero integral vector.

\item The same construction works in every Picard component. Each component contains a finite-order line bundle $P$: for a representative $P_0$, a positive power $P_0^m$ lies in the complex vector group $\Pic^0(Y)$, where division by $m$ is possible. Dividing $P_0$ by the resulting line makes its $m$th power trivial. Replacing $M_r$ by $M_r\otimes P$ retains nontrivial unitary restriction on each fibre and on each specified normalization. Indeed an $m$th power eliminates $P$ and reduces a purported trivialization to the preceding irrationality contradiction.

\item Set $L=K_Y\otimes(M_r\otimes P)^{-1}$. On every full fibre $F$ the Serre-dual bundle for degree two is
\[
 (T_Y\otimes L)|_F^*\otimes\omega_F
       =\Omega_Y^1|_F\otimes(M_r\otimes P)|_F,
\]
because $F$ is a Cartier fibre and $\omega_F=K_Y|_F$. At a smooth torus fibre both the conormal line and the intrinsic cotangent bundle are trivial; nontrivial unitary degree-zero lines on the torus have zero cohomology. The ambient space of sections is therefore zero. At a reduced irreducible normal-crossing fibre, Statement \ref{cdp:st:A.4.5} proves the torsion-free cotangent vanishing for every nontrivial unitary line. The normal-crossing kernel in \ref{cdp:st:A.1.2} is the normalization pushforward, and its sections vanish because the pulled-back unitary line is nontrivial. A holomorphic section of a unitary line on a compact connected complex manifold has constant norm by the maximum principle and is either zero or a trivialization. Thus the ambient group vanishes there as well.

\item On a cyclic-cover neighbourhood of a smooth multiple fibre, the covering map to $Y$ is finite etale. The pulled-back line restricts nontrivially to the covering central torus: finite etale pullback is injective on rational degree-one cohomology by its trace, and $r$ is irrational. The scheme-theoretic inverse image of the full fibre is the corresponding thickening of this torus. Each layer in its ideal filtration is the same ambient cotangent restriction tensored with this nontrivial line, since the covering base coordinate trivializes the conormal. Every layer has zero cohomology. The exact filtration sequences therefore give zero sections on the full thickening. Taking finite-group invariants, an exact functor, proves the required ambient vanishing on the original full fibre.

\item Top-degree coherent base change now gives $R^2f_*(T_Y\otimes L)=0$. For sufficiently large integer $n$, Serre vanishing on the base gives
\[
H^1\bigl(\PP^1,R^1f_*(T_Y\otimes L)\otimes\OO(n)\bigr)=0.
\]
Projection formula and the two-column Leray sequence imply
$H^2(Y,T_Y\otimes L\otimes f^*\OO(n))=0$.
This is an actual twist in the same Picard component up to the fixed finite-order class of $f^*\OO(n)$; choosing $n$ divisible by the finite order of its Chern class keeps the component unchanged. Semicontinuity in the explicit Picard family of C.1 makes the degree-two zero locus dense open in that component. Intersect it with the dense open endpoint-vanishing loci from C.1. On this intersection $\chi(T_Y\otimes L)=-h^1\le0$.

\item Statement \ref{cdp:st:A.4.6} gives Euler number zero or two for each reduced irreducible normal-crossing fibre; the positive case has rational normalization. Our normalization hypothesis excludes that case, since rational smooth surfaces have $H^1(-,\Q)=0$. Smooth torus fibres and the smooth finite etale quotients in the second fibre class have Euler number zero. Therefore \SeriesRef{cdp:st:C2a} gives $c_3(Y)=0$. The index formula in C.1 gives $\chi=0$, forcing $h^1=0$ on the same dense open set. This proves all of \eqref{cdp:eq:C.15}.

\end{enumerate}
\end{proof}
